\documentclass[tikz,margin=3mm]{standalone}
\usepackage{xcolor}
\usepackage{amsmath,amssymb}

\usepackage[utf8]{inputenc}
\usepackage[spanish]{babel}
\usepackage[
    letterpaper,
    top=1.5cm,
    bottom=2cm,
    left=2cm,
    right=2cm
]{geometry}

\usepackage{graphicx}
\usepackage{fancyhdr}
\usepackage[table,xcdraw]{xcolor}
\usepackage{tikz}
\usetikzlibrary{shapes.geometric, arrows, positioning, fit, backgrounds, calc, babel}
\usepackage{ifpdf}
\usepackage{blindtext}
\usepackage[cmex10]{amsmath}
\usepackage{amssymb}
\usepackage{amsfonts}
\usepackage{float}
\usepackage{siunitx}
\usepackage{listings}
\usepackage{hyperref}
\usepackage{subcaption}
\usepackage{booktabs}
\usepackage{url}
\usepackage{wrapfig}
\usepackage{titlesec}
\usepackage[export]{adjustbox}
\usepackage{imakeidx}
\usepackage{parskip}
\usepackage[numbers]{natbib}

\usepackage{ragged2e}
\usepackage{setspace}
\usepackage{array}
\usepackage{longtable}

% Definición de un nuevo tipo de columna para evitar "Underfull \hbox" en tablas
\newcolumntype{L}[1]{>{\RaggedRight\arraybackslash}p{#1}}

\hypersetup{colorlinks=false,bookmarksopen=true,linkbordercolor={1 1 1}}

\newcommand{\rfigura}[1]{Fig.\ref{#1}}
\newcommand{\rtabla}[1]{Tabla \ref{#1}}

% Datos del informe
\newcommand{\Curso}{IEE2913 --- Diseño Eléctrico (Capstone)}
\newcommand{\TituloInforme}{Informe de Avance 1 (5\%)}
\newcommand{\NumeroGrupo}{10}
\newcommand{\NombreProyecto}{SupaClock: Wearable Biométrico Modular}
\newcommand{\Integrantes}{Tomás Avendaño, Benjamín Sepúlveda, Pablo Uribe}
\newcommand{\FechaEntrega}{7 de abril de 2026}

\lstset{
  basicstyle=\ttfamily\small,
  breaklines=true,
  breakatwhitespace=true,
  postbreak=\mbox{\textcolor{red}{$\hookrightarrow$}\space},
  captionpos=b,
  frame=single,
  numbers=left,
  numberstyle=\tiny\color{gray},
  language=C,
  keywordstyle=\color{blue}\bfseries,
  commentstyle=\color{gray}\itshape,
  stringstyle=\color{red!70!black}
}


\begin{document}
\begin{tikzpicture}[
                node distance=1.0cm and 1.5cm,
                block/.style={rectangle, rounded corners, draw=black, thick, fill=yellow!15, text width=2.8cm, align=center, minimum height=1.2cm, font=\footnotesize},
                consumer/.style={rectangle, rounded corners, draw=black, thick, fill=blue!10, text width=2.2cm, align=center, minimum height=0.9cm, font=\footnotesize},
                arrow/.style={thick, ->, >=stealth, red!70},
                dataarrow/.style={thick, ->, >=stealth, blue!60},
                legend/.style={rectangle, draw=black!50, fill=white, inner sep=5pt, font=\scriptsize}
            ]%
            % Fuente
            \node[block, fill=gray!15] (usbc) {USB-C\\5V Entrada};
            \node[block, right=1.5cm of usbc] (bms) {BMS TP4056\\+ DW01A};
            \node[block, right=2.5cm of bms] (lipo) {Batería LiPo\\3.7V, 250 mAh};
            \node[block, below=1.2cm of bms] (ldo) {LDO ME6211\\3.3V Salida};
            \node[block, fill=green!12, below=1.2cm of lipo] (fuel) {Fuel Gauge\\MAX17048};
            %
            % Consumidores (dominio 3.3V)
            \node[consumer, below=1.5cm of ldo, xshift=-3.5cm] (c_mcu) {ESP32-C3/S3\\(MCU)};
            \node[consumer, right=0.5cm of c_mcu] (c_lcd) {ST7789\\(Pantalla)};
            \node[consumer, right=0.5cm of c_lcd] (c_ppg) {MAX30102\\(PPG)};
            \node[consumer, right=0.5cm of c_ppg] (c_temp) {MAX30205\\(Temp)};
            \node[consumer, right=0.5cm of c_temp] (c_imu) {BMI160\\(IMU)};
            \node[consumer, right=0.5cm of c_imu] (c_ecg) {AD8232\\(ECG)};
            %
            % Grupo dominio 3.3V
            \begin{scope}[on background layer]
                \node[rectangle, draw=black!40, dashed, rounded corners, inner sep=0.3cm, fill=blue!3,
                      fit=(c_mcu) (c_lcd) (c_ppg) (c_temp) (c_imu) (c_ecg),
                      label={[font=\scriptsize\bfseries, fill=white, inner sep=2pt]below:Dominio Lógico 3.3V}] (domain) {};
            \end{scope}
            %
            % Conexiones de potencia
            \draw[arrow] (usbc) -- node[above, font=\scriptsize] {5V} (bms);
            \draw[arrow, <->] (bms) -- node[above, font=\scriptsize] {Carga / Descarga} (lipo);
            \draw[arrow] (bms.south) -- node[right, font=\scriptsize] {OUT (3.0--4.2V)} (ldo.north);
            \draw[arrow, line width=1.5pt] (ldo.south) -- node[left, font=\scriptsize\bfseries] {3.3V} ++(0,-0.3) coordinate (spine);
            %
            % Fuel Gauge monitoreo
            \draw[dataarrow] (lipo.south) -- node[right, font=\scriptsize] {V\_CELL} (fuel.north);
            \path (c_ecg.east) ++(0.8cm, 0) coordinate (far_right);
            \path (c_mcu.south) ++(0, -1.0cm) coordinate (under_mcu);
            \draw[dataarrow, dashed] (fuel.east) -| (far_right) |- (under_mcu) node[pos=0.75, below, font=\scriptsize] {I2C (SoC\%)} -| (c_mcu.south);
            %
            % LDO a consumidores
            \foreach \node in {c_mcu, c_lcd, c_ppg, c_temp, c_imu, c_ecg} {
                \draw[arrow, red!40, thin] (spine) -| (\node.north);
            }
            %
            % Leyenda
            \node[legend, below=1.2cm of domain] (leg) {
                \begin{tabular}{cl}
                    \tikz\draw[red!70, thick, ->] (0,0.1) -- (0.4,0.1);   & Flujo de potencia \\
                    \tikz\draw[blue!60, thick, ->] (0,0.1) -- (0.4,0.1);  & Señal de datos    \\
                \end{tabular}
            };
        \end{tikzpicture}
\end{document}
